\chapter{Ruang Afin}\label{ch:b2:affine}

Ruang vektor punya satu titik istimewa --- yakni titik asalnya --- yang
tak diinginkan geometri. Sebuah \emph{\omterm{def:b2:affine:def}{ruang afin}} adalah ruang vektor
yang telah melupakan titik asalnya: sebab titik dan vektor menjadi jenis
yang berbeda, yang dikaitkan penggeseran. Bab pendek ini membangun
kamusnya (yakni titik, \omterm{def:b2:affine:barycenter}{barisentrum}, subruang dan \omterm{def:b2:affine:subspace}{pemetaan afin}), pandangan
\omterm{def:b2:affine:subspace}{afin} atas kecembungan, dan perkakas penggolongan yang dipakai pada
bab geometri berikutnya.

\section{Titik dan vektor}

\begin{definition}\label{def:b2:affine:def}
Sebuah \emph{ruang afin}\index{ruang afin} yang diarahkan ruang
vektor real $E$ adalah himpunan tak kosong $\mathcal{E}$ dengan pemetaan
$(A, B) \mapsto \vect{AB} \in E$ yang memenuhi
\[
\vect{AB} + \vect{BC} = \vect{AC}
\quad \text{(Chasles)},
\qquad
\text{untuk tiap } A,\ B \mapsto \vect{AB} \text{ merupakan bijeksi }
\mathcal{E} \to E .
\]
Kita menulis $B = A + u$ bagi titik tunggal dengan $\vect{AB} = u$.
Adapun \emph{dimensi} $\mathcal{E}$ adalah $\dim E$. Setiap ruang
vektor merupakan ruang \omterm{def:b2:affine:subspace}{afin} atas dirinya sendiri (dengan $\vect{AB} = B - A$); dan setiap
pemilihan titik asal $O \in \mathcal{E}$ mengenali $\mathcal{E}$ dengan
$E$ lewat $M \mapsto \vect{OM}$.
\end{definition}

\begin{example}[Sebuah ruang afin tanpa titik asal
yang alami]\label{ex:b2:affine:planeexample}
Bidang penyelesaian $\mathcal E = \{(x, y, z) \in \R^3 : x + y
+ z = 1\}$ bukan subruang vektor (sebab $0 \notin \mathcal E$),
tetapi ia \omterm{def:b2:affine:def}{ruang afin} yang diarahkan $E = \{x + y + z =
0\}$: sebab untuk $A, B \in \mathcal E$ selisih $\vect{AB} = B
- A$ mendarat di $E$ (karena jumlahnya saling meniadakan), Chasles diwarisi
dari $\R^3$, dan $B \mapsto \vect{AB}$ bijektif ke $E$.
Tak ada titik $\mathcal E$ yang diistimewakan --- sebab sembarang pemilihan
``titik asal'' $O \in \mathcal E$ sama saja berhasil, dan semua
pengenalan $M \mapsto \vect{OM}$ hanya berselisih penggeseran.
Inilah keadaan yang khas: himpunan penyelesaian masalah
linear yang tak homogen (yakni sistem linear, persamaan \omterm{def:b2:diffcalc:differential}{diferensial}
linear pada \cref{ch:b2:diffeq}) bersifat \omterm{def:b2:affine:subspace}{afin}, tak pernah linear,
dan slogan ``penyelesaian khusus ditambah kernel'' persis merupakan
pernyataan $\mathcal F = A + F$ pada definisi berikutnya.
\end{example}

\begin{definition}[Barisentrum]\label{def:b2:affine:barycenter}
Misalkan $(A_i, \lambda_i)_{i \leq k}$ titik terbobot dengan
$\sum\lambda_i \neq 0$. Adapun \emph{barisentrum}\index{barisentrum} $G
= \operatorname{bar}\bigl((A_i, \lambda_i)\bigr)$ adalah titik tunggal
dengan
\[
\sum_i \lambda_i\, \vect{GA_i} = 0
\qquad\text{setara dengan}\qquad
\vect{OG} = \frac{1}{\sum\lambda_i}\sum_i \lambda_i\,\vect{OA_i}
\quad (\text{sembarang } O).
\]
Barisentrum bersifat \emph{asosiatif} (sebab subkelompok titiknya boleh
diganti barisentrum parsialnya dengan bobot yang dijumlahkan) dan
invarian terhadap penskalaan ulang semua bobotnya.
\end{definition}

\begin{proof}[Bukti keberadaan dan rumusnya]
Tetapkan $O$ lalu tulis $s = \sum_i\lambda_i \neq 0$. Menurut Chasles,
\[
\sum_i\lambda_i\,\vect{GA_i} = 0
\iff \sum_i\lambda_i\bigl(\vect{GO} + \vect{OA_i}\bigr) = 0
\iff s\,\vect{OG} = \sum_i\lambda_i\,\vect{OA_i},
\]
yang menentukan $G = O + \frac1s\sum_i\lambda_i\vect{OA_i}$
secara tunggal. \emph{Ketakbergantungannya pada $O$:} sebab untuk titik asal lain
$O'$,
\[
\frac1s\sum_i\lambda_i\,\vect{O'A_i}
= \frac1s\sum_i\lambda_i\bigl(\vect{O'O} + \vect{OA_i}\bigr)
= \vect{O'O} + \frac1s\sum_i\lambda_i\,\vect{OA_i}
= \vect{O'G} :
\]
yakni titik $G$ yang sama. \emph{Keasosiatifannya:} pecahlah himpunan indeksnya
menjadi $I \sqcup J$ dengan $s_I = \sum_{i\in I}\lambda_i \neq 0$,
lalu misalkan $G_I$ \omterm{def:b2:affine:barycenter}{barisentrum} $(A_i, \lambda_i)_{i\in
I}$, sehingga $\sum_{i\in I}\lambda_i\vect{OA_i} =
s_I\,\vect{OG_I}$. Maka
\[
s\,\vect{OG}
= \sum_{i\in I}\lambda_i\vect{OA_i} +
\sum_{j\in J}\lambda_j\vect{OA_j}
= s_I\,\vect{OG_I} + \sum_{j\in J}\lambda_j\vect{OA_j} :
\]
jadi $G$ adalah \omterm{def:b2:affine:barycenter}{barisentrum} $(G_I, s_I)$ bersama $(A_j,
\lambda_j)_{j\in J}$, seperti diklaim. \emph{Penskalaan ulangnya:} mengganti
tiap $\lambda_i$ dengan $t\lambda_i$ (dengan $t \neq 0$) mengalikan $s$
dan jumlah terbobotnya dengan $t$, jadi $\vect{OG}$ tak berubah.
\end{proof}

\begin{remark}[Jebakan yang sering muncul]
Dua jebakan mengelilingi definisinya. Pertama, bila bobotnya berjumlah
\emph{nol}, maka tak ada \omterm{def:b2:affine:barycenter}{barisentrum}: sebab pemetaan $O \mapsto
\sum\lambda_i\vect{OA_i}$ lalu tak bergantung pada $O$ dan
mendefinisikan sebuah \emph{vektor}, bukan titik --- misalnya
$(A, -1;\ B, 1)$ menyandikan $\vect{AB}$. Jadi melacak yang mana
dari kedua objek itu yang dihasilkan sebuah perhitungan adalah separuh dari
kebersihan barisentrik. Kedua, bobotnya hanya bermakna hingga
sebuah faktor tak nol bersama; jadi rumus seperti ``koordinat
$G$ adalah $\lambda_1, \dots, \lambda_k$'' mengandaikan sebuah
penormalan (biasanya $\sum\lambda_i = 1$), dan lupa
menormalkannya adalah sumber baku bagi nisbah yang salah pada
sebuah gambar.
\end{remark}

\begin{definition}[Subruang afin; pemetaan afin]\label{def:b2:affine:subspace}
Sebuah \emph{subruang afin} adalah himpunan $\mathcal{F} = A + F = \{A + u :
u \in F\}$ dengan $F$ sebuah subruang vektor (yakni \emph{arahnya});
setara dengan itu, sebuah himpunan tak kosong yang stabil terhadap \omterm{def:b2:affine:barycenter}{barisentrum}. Adapun subruang
afin $\R^n$ persis merupakan himpunan penyelesaian sistem linear
$MX = B$ (lewat Tahun ke-1: penyelesaian khusus ditambah kernel). Sebuah pemetaan $f \colon
\mathcal{E} \to \mathcal{E}'$ disebut \emph{afin}\index{pemetaan afin}
apabila ia memelihara \omterm{def:b2:affine:barycenter}{barisentrum} --- setara dengan apabila
\[
f(A + u) = f(A) + \varphi(u)
\]
bagi sebuah pemetaan linear (yang tunggal) $\varphi = \vec f$, yakni \emph{bagian
linearnya}. Pemetaan afin $\R^n$: yaitu $X \mapsto MX + C$. Komposisinya bersifat
afin dengan bagian linear yang dikomposisikan; dan $f$ bijektif bila dan hanya bila $\vec f$
bijektif.
\end{definition}

\begin{proof}[Bukti kesetaraannya bagi pemetaan]
Bila $f(A + u) = f(A) + \varphi(u)$: maka untuk \omterm{def:b2:affine:barycenter}{barisentrum} $G$ milik $(A_i,
\lambda_i)$, dengan menguraikan setiap titiknya dari $A$, $f(G) = f(A) +
\varphi(\vect{AG})$ dan $\varphi(\vect{AG}) =
\frac{\sum\lambda_i\varphi(\vect{AA_i})}{\sum\lambda_i}$: jadi $f(G)$ adalah
\omterm{def:b2:affine:barycenter}{barisentrum} petanya. Sebaliknya, tetapkan $A$ lalu definisikan
$\varphi(u) = \vect{f(A)\,f(A + u)}$. \emph{Kehomogenannya:} sebab $A + tu =
\operatorname{bar}\bigl(A, 1-t;\ A + u, t\bigr)$ untuk setiap $t$ yang
real, jadi pemeliharaan \omterm{def:b2:affine:barycenter}{barisentrumnya} (dengan bobot real sembarang,
sebagaimana dihipotesiskan) memberi $\varphi(tu) = t\,\varphi(u)$ secara langsung.
\emph{Keaditifannya:} sebab $A + u + v = \operatorname{bar}\bigl(A + 2u,
\tfrac12;\ A + 2v, \tfrac12\bigr)$, jadi $\varphi(u + v) =
\tfrac12\varphi(2u) + \tfrac12\varphi(2v) = \varphi(u) +
\varphi(v)$, dengan memakai kehomogenannya. Karena itu $\varphi$ bersifat linear.
\end{proof}

\begin{remark}
Buktinya memakai \omterm{def:b2:affine:barycenter}{barisentrum} dengan bobot \emph{real yang
sembarang}: sebab langkah kehomogenannya mengambil $t$ di luar $\intcc01$.
Bila sebuah pemetaan hanya diandaikan memelihara \omterm{def:b2:affine:barycenter}{barisentrum} berbobot
taknegatif --- yakni setara dengan memelihara titik tengah dan ruas
--- maka kelinearan pemetaan vektornya tak lagi gratis:
sebab kita hanya memperoleh kelinearan atas $\Q$, dan sebuah hipotesis kekontinuan
diperlukan untuk menyimpulkannya, persis seperti pada
\cref{exo:b2:affine:5}. Jadi membedakan ``memelihara semua
\omterm{def:b2:affine:barycenter}{barisentrum}'' dari ``memelihara kombinasi \omterm{def:b2:affine:convex}{cembung}'' adalah
kehalusan kecil tetapi nyata pada kosakata \omterm{def:b2:affine:subspace}{afinnya}.
\end{remark}

\begin{example}[Geometri barisentrum yang klasik]\label{ex:b2:affine:median}
Titik berat segitiga $ABC$ adalah \omterm{def:b2:affine:barycenter}{barisentrum} $G =
\operatorname{bar}(A,1; B,1; C,1)$. Keasosiatifannya bersama
titik tengah $A' = \operatorname{bar}(B, 1; C, 1)$ menunjukkan
\[
G = \operatorname{bar}(A, 1;\ A', 2) :
\]
jadi $G$ terletak pada garis berat $AA'$ pada dua pertiganya --- dan demikian pula
bagi dua garis berat lainnya: sehingga ketiga garis beratnya berkonkuren, dalam
satu baris kalkulus barisentrik.
\end{example}

\begin{example}[Bimedian sebuah
segi empat]\label{ex:b2:affine:bimedians}
Misalkan $ABCD$ sembarang segi empat (bidang atau bukan!) lalu tinjaulah
\emph{bimedian}nya: yakni ruas yang menghubungkan titik tengah sisi
berhadapannya, $M_{AB}M_{CD}$ dan $M_{BC}M_{DA}$. Perkenalkanlah
\omterm{def:b2:affine:barycenter}{barisentrum} $G$ milik $(A,1; B,1; C,1; D,1)$ lalu kelompokkanlah
bobotnya lewat dua cara:
\[
G = \operatorname{bar}\bigl(M_{AB}, 2;\ M_{CD}, 2\bigr)
= \operatorname{bar}\bigl(M_{BC}, 2;\ M_{DA}, 2\bigr) :
\]
jadi $G$ adalah titik tengah \emph{kedua} bimediannya --- sehingga kedua
bimediannya selalu saling membagi dua, dan segi empat yang dibentuk
keempat titik tengahnya adalah jajaran genjang (yang diagonalnya adalah
bimediannya). Jadi tanpa analisis kasus, tanpa koordinat, dan
hujahnya bertahan tanpa berubah bagi segi empat yang bersilang di
$\R^3$, tempat bukti berbasis gambar sudah akan
peka: sebab keasosiatifannya tak peduli pada dimensinya.
\end{example}

\begin{example}[Menggolongkan sebuah pemetaan afin, dari awal sampai
akhir]\label{ex:b2:affine:dilation}
Misalkan $f(x, y) = (2x - 1,\ 3y - 4)$ pada $\R^2$. Bagian linearnya
adalah $\varphi = \operatorname{diag}(2, 3)$, yang \omterm{def:b2:reduction:eigen}{spektrumnya}
$\{2, 3\}$ menghindari $1$: jadi menurut kriteria titik tetap yang dibuktikan
di bawah (\cref{prop:b2:affine:fixedpoint}), $f$ punya tepat satu
titik tetap, yang ditemukan dengan menyelesaikan
\[
x = 2x - 1, \qquad y = 3y - 4
\qquad\Longrightarrow\qquad \Omega = (1, 2).
\]
Setelah dipusatkan ulang di $\Omega$ (tetapkan $x = 1 + u$, $y = 2 + v$):
\[
f(1 + u,\ 2 + v) = (1 + 2u,\ 2 + 3v) :
\]
dalam kerangka di $\Omega$, $f$ \emph{adalah} bagian linearnya, yakni sebuah
pemuaian anisotropik yang meregang $2$ kali mendatar dan $3$ kali
tegak dari pusat $(1, 2)$. Pelajaran umumnya: sebuah
\omterm{def:b2:affine:subspace}{pemetaan afin} adalah ``pemetaan linear ditambah data lokasi'', dan
data lokasinya runtuh menjadi satu titik asal yang terpilih baik setiap kali
$1$ bukan \omterm{def:b2:reduction:eigen}{nilai eigennya}. Sebaliknya, menggeser titik asalnya
secara buruk justru \emph{menciptakan} suku tetapnya: jadi geometri \omterm{def:b2:affine:subspace}{afin} adalah
seni memilih tempat menaruh $0$.
\end{example}

\begin{remark}[Metode: kekonkurenan dan kesegarisan lewat
barisentrum]
Contoh \cref{ex:b2:affine:median} adalah sebuah wujud resep
umum. Untuk membuktikan bahwa tiga cevian sebuah segitiga
berkonkuren, pamerkanlah \emph{satu} sistem terbobot $(A,
\alpha; B, \beta; C, \gamma)$ lalu pakailah keasosiatifannya lewat tiga
cara: sebab mengelompokkan $(B, C)$ menunjukkan \omterm{def:b2:affine:barycenter}{barisentrumnya} terletak pada
cevian dari $A$, mengelompokkan $(C, A)$ pada cevian dari $B$,
dan mengelompokkan $(A, B)$ pada yang ketiga. Untuk garis beratnya, sistem
$(A, 1; B, 1; C, 1)$ mengerjakan semuanya; sedangkan untuk cevian yang memotong
sisinya dengan nisbah yang ditentukan, bobotnya dibaca dari
nisbahnya. Untuk membuktikan tiga titik \emph{segaris}, tulislah satu sebagai
\omterm{def:b2:affine:barycenter}{barisentrum} dua lainnya (\cref{exo:b2:affine:2}), atau pakailah
kriteria \omterm{def:b2:linalg:det}{determinan} pada \cref{exo:b2:affine:11}. Kedua
resepnya mengganti kecerdikan geometri dengan tata buku bobot ---
dan persis untuk itulah kalkulus barisentrik ada.
\end{remark}

\section{Kecembungan, secara afin}

\begin{definition}\label{def:b2:affine:convex}
Sebuah himpunan bagian $C$ pada \omterm{def:b2:affine:def}{ruang afin} disebut \emph{cembung}\index{himpunan cembung} apabila ia memuat
setiap \omterm{def:b2:affine:barycenter}{barisentrum} berbobot \emph{taknegatif} atas titiknya ---
setara dengan itu, setiap ruas $\intcc{A}{B} = \{\operatorname{bar}(A,
1-t; B, t) : t \in \intcc{0}{1}\}$ di antara titiknya. Adapun
\emph{bungkus cembung}\index{bungkus cembung} $\operatorname{conv}(S)$ adalah himpunan semua
\omterm{def:b2:affine:barycenter}{barisentrum} berbobot taknegatif atas titik $S$ --- yakni himpunan
cembung terkecil yang memuat $S$.
\end{definition}

\begin{example}[Epigraf itu himpunan
cembung]\label{ex:b2:affine:epigraph}
Daerah $C = \{(x, y) : y \geq x^2\}$ di atas parabolanya
bersifat \omterm{def:b2:affine:convex}{cembung}: sebab untuk $(x_1, y_1), (x_2, y_2) \in C$ dan $t \in
\intcc01$, ketaksamaan kecembungan fungsi kuadratnya
memberi
\[
\bigl((1-t)x_1 + tx_2\bigr)^2 \leq (1-t)x_1^2 + tx_2^2
\leq (1-t)y_1 + ty_2 ,
\]
jadi \omterm{def:b2:affine:barycenter}{barisentrumnya} tetap di atas parabolanya. Adapun perhitungannya
bersifat umum: $\{y \geq f(x)\}$ \omterm{def:b2:affine:convex}{cembung} persis ketika $f$ merupakan
fungsi \omterm{def:b2:affine:convex}{cembung} --- jadi \emph{himpunan} \omterm{def:b2:affine:convex}{cembung} dan \emph{fungsi}
\omterm{def:b2:affine:convex}{cembung} (\cref{ch:b2:realfun}) adalah dua wajah satu
gagasan, dengan epigrafnya sebagai kamusnya. Inilah alasan geometris
mengapa garis penyangga ada bagi fungsi \omterm{def:b2:affine:convex}{cembung}, yakni fakta
yang kelak membuktikan ketaksamaan Jensen pada
\cref{ch:b2:randomvar}.
\end{example}

\begin{example}[Pembangun yang mubazir bagi sebuah bungkus
cembung]\label{ex:b2:affine:squarehull}
Misalkan $S = \{(0,0), (2,0), (2,2), (0,2), (1,1)\}$. Titik
kelimanya adalah \omterm{def:b2:affine:barycenter}{barisentrum}
\[
(1,1) = \operatorname{bar}\bigl((0,0), \tfrac12;\ (2,2),
\tfrac12\bigr),
\]
jadi ia sudah terletak di bungkus keempat lainnya:
sehingga $\operatorname{conv}(S)$ adalah bujur sangkar dengan keempat pojoknya
sebagai titik sudutnya. Secara umum, titik $S$ yang merupakan
\omterm{def:b2:affine:barycenter}{barisentrum} berbobot taknegatif atas titik \emph{lain} milik
$S$ dapat dihapus tanpa mengubah bungkusnya; sedangkan titik yang
tak pernah dapat dihapus (di sini keempat pojoknya) adalah
\emph{titik ekstrem} bungkusnya. Menentukannya adalah perhitungan
\omterm{def:b2:affine:barycenter}{barisentrum} yang murni: misalnya $(2,0)$ tak dapat dituliskan sebagai
$\operatorname{bar}$ titik sisanya dengan bobot yang
taknegatif, karena koordinat pertamanya akan memaksa seluruh bobotnya
ke titik yang $x = 2$, lalu koordinat keduanya
gagal. Jadi pertanyaan kecembungan menyusut, berulang kali, menjadi
menyelesaikan sistem terbobot yang kecil.
\end{example}

\begin{theorem}[Carathéodory]\label{thm:b2:affine:caratheodory}
Pada \omterm{def:b2:affine:def}{ruang afin} berdimensi $n$, setiap titik
$\operatorname{conv}(S)$ adalah \omterm{def:b2:affine:barycenter}{barisentrum} atas paling banyak $n + 1$ titik
$S$.\index{teorema Carathéodory}
\end{theorem}

\begin{proof}
Misalkan $G = \operatorname{bar}(A_0, \lambda_0; \dots; A_k, \lambda_k)$
dengan $\lambda_i > 0$, $\sum\lambda_i = 1$, dan $k + 1 > n + 1$
titik. Adapun $k$ vektor $\vect{A_0A_i}$ (untuk $i \geq 1$) bergantungan
(sebab $k > n$): jadi $\sum_{i\geq1}\mu_i \vect{A_0A_i} = 0$ secara tak trivial;
lalu dengan menetapkan $\mu_0 = -\sum_{i\geq1}\mu_i$, kita memperoleh bobot $(\mu_i)$
dengan $\sum\mu_i = 0$ dan $\sum \mu_i\,\vect{OA_i} = 0$ (untuk sembarang $O$), yang tak
semuanya nol. Maka untuk setiap $t$ yang real bobot $\lambda_i - t\mu_i$
tetap berjumlah $1$ dan, karena $\sum_i\mu_i\vect{OA_i} = 0$,
\[
\sum_i(\lambda_i - t\mu_i)\,\vect{OA_i}
= \sum_i\lambda_i\,\vect{OA_i} :
\]
jadi keduanya menghasilkan titik $G$ yang \emph{sama}. Kini geserlah $t$ dari
$0$: suatu $\mu_i$ bersifat positif (sebab jumlahnya nol dan tak
semuanya nol), jadi
\[
t^* = \min\Bigl\{\frac{\lambda_i}{\mu_i} : \mu_i > 0\Bigr\}
\]
terdefinisi baik dan positif. Di $t = t^*$: untuk indeks dengan
$\mu_i > 0$ berlaku $\lambda_i - t^*\mu_i \geq 0$ berkat keminimalannya,
dengan kesamaan pada indeks yang meminimumkannya; sedangkan untuk indeks dengan $\mu_i
\leq 0$ berlaku $\lambda_i - t^*\mu_i \geq \lambda_i > 0$. Jadi semua
bobotnya tetap taknegatif dan sekurangnya satu telah mati: sehingga $G$
dituliskan ulang sebagai \omterm{def:b2:affine:barycenter}{barisentrum} atas titik yang lebih sedikit. Ulangilah selama lebih
dari $n + 1$ titik tersisa.
\end{proof}

\begin{example}
Pada bidangnya (dengan $n = 2$): setiap titik \omterm{def:b2:affine:convex}{bungkus cembung} sebuah himpunan
berhingga terletak di sebuah segitiga yang titik sudutnya di himpunan itu --- yakni isi
geometris Carathéodory, yang dipakai dalam pengoptimalan dan peluang
(yakni campuran) sekaligus.
\end{example}

\begin{example}[Menjalankan algoritma
Carathéodory]\label{ex:b2:affine:caratheodoryrun}
Tulislah pusat bujur sangkar pada
\cref{ex:b2:affine:squarehull} dengan keempat pojoknya $A_1 =
(0,0)$, $A_2 = (2,0)$, $A_3 = (2,2)$, $A_4 = (0,2)$:
\[
(1,1) = \operatorname{bar}\bigl(A_1, \tfrac14;\ A_2,
\tfrac14;\ A_3, \tfrac14;\ A_4, \tfrac14\bigr),
\]
yakni empat titik pada dimensi $2$ --- satu terlalu banyak. Adapun resep
buktinya meminta bobot $(\mu_i)$ dengan $\sum\mu_i = 0$ dan
$\sum\mu_i\vect{OA_i} = 0$: di sini $\mu = (1, -1, 1, -1)$ berhasil
(sebab kedua diagonalnya berbagi titik tengahnya). Menggeser $\lambda_i
\mapsto \lambda_i - t\mu_i$ menjaga \omterm{def:b2:affine:barycenter}{barisentrumnya} tetap untuk
setiap $t$; lalu nilai terterima yang ekstrem $t = \frac14$ membuat
bobotnya $(0, \tfrac12, 0, \tfrac12)$, yang membunuh $A_1$ dan
$A_3$ sekaligus:
\[
(1,1) = \operatorname{bar}\bigl(A_2, \tfrac12;\ A_4,
\tfrac12\bigr),
\]
yakni penyajian lewat dua titik --- bahkan lebih baik daripada tiga
yang dijamin teoremanya, karena pusatnya kebetulan
terletak pada sebuah ruas di antara pembangunnya. Adapun algoritmanya
sepenuhnya mekanis: carilah sebuah kebergantungan, geserlah sampai sebuah bobot
mati, lalu ulangi.
\end{example}

\section{Perkakas penggolongan afin}

\begin{proposition}[Titik tetap pemetaan afin]\label{prop:b2:affine:fixedpoint}
Misalkan $f$ endomorfisma \omterm{def:b2:affine:subspace}{afin} sebuah \omterm{def:b2:affine:def}{ruang afin} berdimensi
hingga dengan bagian linear $\varphi$. Bila $1 \notin
\operatorname{Sp}(\varphi)$, maka $f$ punya tepat satu titik tetap
$\Omega$, dan dalam pemvektoran di $\Omega$, $f$ \emph{adalah}
bagian linearnya. (Adapun penggeseran, dengan $\varphi = \mathrm{id}$ dan tanpa
titik tetap, merupakan halangan dasarnya.)
\end{proposition}

\begin{proof}
Tetapkan $O$ lalu tulis $f(O + x) = f(O) + \varphi(x)$. Titik
$O + x$ bersifat tetap bila dan hanya bila $O + x = f(O) + \varphi(x)$, yakni
\[
(\mathrm{id} - \varphi)(x) = \vect{O f(O)} .
\]
Pada dimensi hingga, $\mathrm{id} - \varphi$ terbalikkan
bila dan hanya bila $0$ bukan \omterm{def:b2:reduction:eigen}{nilai eigen} $\mathrm{id} - \varphi$, yakni bila
$1 \notin \operatorname{Sp}\varphi$ --- dan pada kasus itu
persamaan pada displainya punya tepat satu penyelesaian $x^*$, yang memberi
titik tetap tunggal $\Omega = O + x^*$. Adapun pemusatan ulangnya: untuk sembarang
vektor $u$,
\[
f(\Omega + u) = f(\Omega) + \varphi(u) = \Omega + \varphi(u),
\]
jadi dalam kerangka bertitik asal $\Omega$ pemetaannya berbunyi $u \mapsto
\varphi(u)$: yakni linear murni. Sedangkan ketika $1 \in
\operatorname{Sp}\varphi$, entah tak ada titik tetap (sebab
persamaan pada displainya boleh jadi tak terselesaikan, seperti untuk penggeseran)
atau justru seluruh subruang \omterm{def:b2:affine:subspace}{afin} titik tetapnya ada (tambahkanlah sembarang \omterm{def:b2:reduction:eigen}{vektor eigen}
bernilai eigen $1$ pada sebuah penyelesaian): jadi ketunggalannya persis
merupakan syarat spektralnya.
\end{proof}

\begin{example}[Isometri bidang, dirampungkan]\label{ex:b2:affine:isometries}
Sebuah isometri \omterm{def:b2:affine:subspace}{afin} bidang Euklides punya bagian linear di
$O(2)$: yakni sebuah perputaran $R_\theta$ atau sebuah pencerminan (lewat jilid Tahun ke-1). Bila
$\theta \neq 0$: maka $1 \notin \operatorname{Sp} R_\theta$, jadi pemetaannya
merupakan \emph{perputaran terhadap satu pusat yang tunggal}
(\cref{prop:b2:affine:fixedpoint}). Bila bagian linearnya sebuah
pencerminan: maka entah sebuah pencerminan pada sebuah sumbu (jika titik tetapnya ada) atau
sebuah \emph{pencerminan geser} (yakni pencerminan yang dikomposisikan dengan penggeseran
sepanjang sumbunya, tanpa titik tetap). Bersama penggeserannya, inilah
penggolongan \omterm{def:b2:metric:complete}{lengkap} isometri bidang.
\end{example}

\begin{remark}[Isometri bidangnya, sekilas]
Mengumpulkan kasusnya: identitas; penggeseran (dengan $\vec f =
\mathrm{id}$, tanpa titik tetap kecuali yang trivial); perputaran
(dengan bagian linear $R_\theta$ dan $\theta \neq 0$: yakni satu pusat);
pencerminan (dengan bagian linear sebuah pencerminan, dan sebuah garis titik
tetap); pencerminan geser (dengan bagian linear yang sama, tanpa titik
tetap). Jadi empat keluarga ditambah identitasnya, yang masing-masing dikenali lewat
dua data belaka: yakni bagian linearnya dan himpunan titik tetapnya ---
inilah pola \cref{prop:b2:affine:fixedpoint} yang dibuat
menyeluruh.
\end{remark}

\begin{example}[Sebuah pencerminan geser, tertangkap
basah]\label{ex:b2:affine:glide}
Misalkan $f(x, y) = (y + 1,\ x + 1)$. Bagian linearnya $(x, y)
\mapsto (y, x)$ adalah pencerminan pada diagonal $y = x$, jadi
$1 \in \operatorname{Sp}\vec f$ dan
\cref{prop:b2:affine:fixedpoint} membisu. Titik tetapnya
akan menuntut $x = y + 1$ dan $y = x + 1$ sekaligus:
yang mustahil --- jadi tak ada, sehingga $f$ bukan pencerminan. Adapun mengkuadratkannya
menyelesaikan penggolongannya:
\[
f\bigl(f(x, y)\bigr) = f(y + 1,\ x + 1) = (x + 2,\ y + 2),
\]
yakni penggeseran sebesar $(2, 2)$: jadi $f$ adalah \emph{pencerminan
geser} yang bersumbu garis $y = x$ (yang digeser
sepantasnya: sebab titik tengah $M$ dan $f(M)$ selalu terletak pada
$y = x + {}$konstanta, di sini $y = x$, seperti diperiksa pada $M = (0,
0) \mapsto (1,1)$) dengan vektor geser $(1, 1)$, yakni separuh $f
\circ f$. Bandingkan dengan \cref{exo:b2:affine:6}, tempat bagian
linear yang sama tetapi konstanta yang berbeda menghasilkan pencerminan
yang jujur: jadi dengan hadirnya \omterm{def:b2:reduction:eigen}{nilai eigen} $1$, suku tetapnya
memutuskan segalanya.
\end{example}

\begin{example}[Rekursi afin itu dinamika
afin]\label{ex:b2:affine:recursion}
Rekursi klasik $u_{n+1} = au_n + b$ (dengan $a \neq 1$)
mengiterasikan \omterm{def:b2:affine:subspace}{pemetaan afin} $f(x) = ax + b$ pada garisnya, yang
bagian linearnya $a$ menghindari \omterm{def:b2:reduction:eigen}{nilai eigen} $1$: jadi ada satu
titik tetap tunggal $\omega = \frac{b}{1-a}$, lalu memusatkan ulang di sana
(yakni kasus berdimensi satu proposisi di atas) mengubah
$f$ menjadi perkalian dengan $a$:
\[
u_{n+1} - \omega = a\,(u_n - \omega)
\qquad\Longrightarrow\qquad
u_n = \omega + a^n(u_0 - \omega) .
\]
Untuk $u_{n+1} = \frac{u_n}2 + 3$: berlaku $\omega = 6$ dan $u_n = 6 +
(u_0 - 6)2^{-n} \to 6$. Adapun resep yang diajarkan bagi rekursi
semacam itu pada \cref{ch:b2:series} --- yakni ``kurangkan titik
tetapnya'' --- persis merupakan pemvektoran sebuah \omterm{def:b2:affine:subspace}{pemetaan afin}
di titik tetapnya; sedangkan kekonvergenannya untuk $\abs a < 1$ adalah
gejala kontraksi yang diubah \cref{ch:b2:metric} menjadi
teorema titik tetap Banach. Jadi satu gagasan, tiga bab.
\end{example}

\begin{example}[Menemukan pusat sebuah perputaran]\label{ex:b2:affine:rotationcenter}
Misalkan $f(x, y) = (-y + 2,\ x)$. Bagian linearnya adalah
$\varphi(x, y) = (-y, x)$: yakni perputaran bersudut $\frac\pi2$,
yang \omterm{def:b2:reduction:eigen}{spektrumnya} $\{\iu, -\iu\}$ menghindari $1$. Menurut
\cref{prop:b2:affine:fixedpoint} ada tepat satu titik
tetap: sebab $x = -y + 2$ dan $y = x$ memberi $x = 1$ dan $y = 1$, jadi
$\Omega = (1, 1)$, dan $f$ \emph{adalah} perputaran berpusat
$(1, 1)$ dan bersudut $\frac\pi2$. Pelajaran umumnya: ketika $1
\notin \operatorname{Sp}\vec f$, menggolongkan $f$ berharga satu
sistem linear --- sebab geometrinya seluruhnya ada pada bagian linearnya,
dan aritmetikanya seluruhnya pada penentuan pusatnya.
\end{example}

\begin{remark}[Di mana bahasa afin dipakai berikutnya]
\omterm{def:b2:affine:barycenter}{Barisentrum} dan \omterm{def:b2:affine:subspace}{pemetaan afin} adalah tata bahasa bab
geometri berikutnya: sebab garis dan bidang singgung adalah objek \omterm{def:b2:affine:subspace}{afin}
(\cref{ch:b2:curves,ch:b2:surfaces}), sebuah penggantian peubah
\omterm{def:b2:affine:subspace}{afin} mengalikan luas dan volumenya dengan
$\abs{\det \vec f}$ (\cref{ch:b2:multint}), dan nilai harapan
adalah \omterm{def:b2:affine:barycenter}{barisentrum} dengan bobot yang diberikan sebuah hukum peluang, dan itulah
sebabnya kecembungan mengatur ketaksamaan Jensen
(\cref{ch:b2:randomvar}). Adapun pada jilid Tahun ke-3 kosakata
kecembungan yang sama mengusung kajian \omterm{def:b2:nvs:norm}{norma} $L^p$ dan
ketaksamaan integral.
\end{remark}

\begin{remark}[Pandangan ke depan di dalam jilid ini]
Dua benang meninggalkan bab ini. Benang \emph{\omterm{def:b2:affine:subspace}{afinnya}}: yakni
garis singgung (\cref{ch:b2:curves}) dan bidang singgung
(\cref{ch:b2:surfaces}) merupakan subruang \omterm{def:b2:affine:subspace}{afin} yang dilekatkan pada
objek nonlinear, dan penggolongan kuadrik pada bab
permukaan berjalan atas persamaan pusat bab ini
$A\Omega = -b$. Adapun benang \emph{\omterm{def:b2:affine:convex}{cembungnya}} lebih panjang: sebab kecembungan
setengah bidang dan cakram menggerakkan teori Helly pada
soal akhir pekan; kecembungan fungsi memberi ketaksamaan Jensen
(\cref{ch:b2:randomvar}); dan teorema terakhir
buku ini --- yakni kriteria kepunahan bagi proses bercabang
(\cref{ch:b2:genfun}) --- diputuskan oleh kedudukan sebuah
kurva \omterm{def:b2:affine:convex}{cembung} terhadap diagonalnya, yakni gambaran yang termasuk
bab ini sebanyak ia termasuk peluang. \omterm{def:b2:affine:barycenter}{Barisentrumnya}
kembali di sana pula: sebab nilai harapan adalah \omterm{def:b2:affine:barycenter}{barisentrum} dengan
bobot peluang.
\end{remark}

\section{Latihan}

\begin{exercise}[$\star$]\label{exo:b2:affine:1}
Pada $\R^3$, apakah yang berikut merupakan subruang \omterm{def:b2:affine:subspace}{afin}? Berikanlah arah dan
dimensinya. $\{x + y + z = 1\}$; $\;\{x + y + z = 1,\ x - z =
3\}$; $\;\{x^2 + y^2 = 1\}$; dan himpunan penyelesaian $MX = B$ bagi
sebuah sistem yang serasi.
\end{exercise}

\begin{exercise}[$\star$]\label{exo:b2:affine:2}
Buktikan bahwa tiga titik berbeda $A, B, C$ pada sebuah \omterm{def:b2:affine:def}{ruang afin}
bersifat segaris bila dan hanya bila $C$ merupakan \omterm{def:b2:affine:barycenter}{barisentrum} $A$ dan $B$, yakni bila vektor
$\vect{AB}, \vect{AC}$ bergantungan. Lalu turunkan tata buku bobot
bergaya Menelaus: bila $C = \operatorname{bar}(A, 1 - t; B, t)$, tentukanlah letak
$C$ untuk $t = \frac12$, $t = 2$, dan $t = -1$.
\end{exercise}

\begin{exercise}[$\star$]\label{exo:b2:affine:3}
Misalkan $f$ \omterm{def:b2:affine:subspace}{pemetaan afin} $\R^2$ yang diberikan $f(X) = MX + C$ dengan
$M = \frac12\begin{pmatrix} 1 & 1\\ 1 & 1\end{pmatrix}$ dan $C =
(1, 0)^{\mathsf T}$. Tentukan peta $f$, titik tetapnya
(bila ada), dan $f \circ f$.
\end{exercise}

\begin{exercise}[$\star\star$]\label{exo:b2:affine:4}
(Keasosiatifan beraksi) Pada segitiga $ABC$, misalkan $I, J, K$ membagi
$BC$, $CA$, $AB$ dengan nisbah $\vect{BI} = \frac13\vect{BC}$,
$\vect{CJ} = \frac13\vect{CA}$, dan $\vect{AK} = \frac13\vect{AB}$.
Nyatakanlah $I, J, K$ sebagai \omterm{def:b2:affine:barycenter}{barisentrum} lalu hitunglah \omterm{def:b2:affine:barycenter}{barisentrum}
$(I,1;J,1;K,1)$: apa yang kautemukan, dan mengapa itu dapat diramalkan?
\end{exercise}

\begin{exercise}[$\star\star$]\label{exo:b2:affine:5}
Buktikan bahwa pemetaan $f \colon \R^n \to \R^n$ yang memelihara \emph{titik tengah}
(yakni $f\bigl(\frac{A+B}{2}\bigr) = \frac{f(A) + f(B)}{2}$) dan
bersifat \emph{\omterm{def:b2:metric:continuity}{kontinu}} pastilah \omterm{def:b2:affine:subspace}{afin}. \emph{(Tunjukkan bahwa pemetaan vektornya $u \mapsto
f(O + u) - f(O)$ bersifat aditif lewat titik tengahnya, lalu
homogen atas $\Q$, lalu homogen atas $\R$ berkat kekontinuannya ---
yakni strategi kepadatan yang sama seperti untuk persamaan fungsional Cauchy pada
jilid Tahun ke-1; turunkanlah ulang langkah yang diperlukan di sini.)}
\end{exercise}

\begin{exercise}[$\star\star$]\label{exo:b2:affine:6}
Golongkanlah \omterm{def:b2:affine:subspace}{pemetaan afin} $f(x, y) = (y + 1,\; x - 1)$ pada bidang
Euklides: yakni bagian linearnya, titik tetapnya, dan sifat geometrisnya
(pencerminan? geser?). Hitunglah $f \circ f$ lalu simpulkan.
\end{exercise}

\begin{exercise}[$\star\star\star$]\label{exo:b2:affine:7}
(Radon) Misalkan $A_1, \dots, A_{n+2}$ titik pada sebuah \omterm{def:b2:affine:def}{ruang afin}
berdimensi $n$. Buktikan bahwa keduanya dapat dipecah menjadi dua kelompok
yang saling lepas dan \omterm{def:b2:affine:convex}{bungkus cembungnya} berpotongan. \emph{(Seperti pada bukti
Carathéodory, carilah bobot $\mu_i$ yang tak semuanya nol, dengan $\sum\mu_i = 0$ dan
$\sum\mu_i\vect{OA_i} = 0$; lalu pisahkanlah bobot positif dan negatifnya
kemudian normalkanlah kedua ruasnya.)}
\end{exercise}

\begin{exercise}[$\star\star\star$]\label{exo:b2:affine:8}
Misalkan $f$ endomorfisma \omterm{def:b2:affine:subspace}{afin} $\R^n$ dengan $f \circ f = f$.
Buktikan bahwa $f$ adalah proyeksi \omterm{def:b2:affine:subspace}{afin} ke subruang \omterm{def:b2:affine:subspace}{afin}
$\operatorname{Fix}(f) = \operatorname{im} f$ sepanjang arah
$\ker\vec f$, dan bahwa sebaliknya semua proyeksi semacam itu bersifat
idempoten. \emph{(Tunjukkanlah lebih dulu bahwa $\operatorname{im} f$ terdiri
atas titik tetap.)}
\end{exercise}

\begin{exercise}[$\star$]\label{exo:b2:affine:9}
Misalkan $G = \operatorname{bar}(A, 1;\ B, 2;\ C, 3)$ pada segitiga
$ABC$. Dengan memakai keasosiatifannya, tunjukkanlah bahwa garis $AG$ memotong $BC$
di $M = \operatorname{bar}(B, 2;\ C, 3)$, lalu tentukan letak $G$ pada
ruas $\intcc AM$; tentukan pula letak perpotongan $BG$
dengan $CA$.
\end{exercise}

\begin{exercise}[$\star\star$]\label{exo:b2:affine:10}
Untuk $\lambda \neq 0$, \emph{homoteti} $h_{\Omega,
\lambda}$ adalah \omterm{def:b2:affine:subspace}{pemetaan afin} yang menetapkan $\Omega$ dengan bagian linear
$\lambda\,\mathrm{id}$. Buktikan bahwa komposisi
$h_{\Omega', \mu} \circ h_{\Omega, \lambda}$ merupakan homoteti
bernisbah $\lambda\mu$ ketika $\lambda\mu \neq 1$, dan sebuah penggeseran
ketika $\lambda\mu = 1$; lalu pada kasus $\lambda = \mu = -1$ (yakni dua
pencerminan titik), hitunglah vektor geserannya.
\end{exercise}

\begin{exercise}[$\star\star$]\label{exo:b2:affine:11}
(Menelaus) Pada segitiga $ABC$, misalkan $A' \in (BC)$, $B' \in
(CA)$, $C' \in (AB)$, yang semuanya berbeda dari titik sudutnya, lalu
definisikan $\alpha, \beta, \gamma$ lewat $\vect{A'B} =
\alpha\,\vect{A'C}$, $\vect{B'C} = \beta\,\vect{B'A}$, dan
$\vect{C'A} = \gamma\,\vect{C'B}$. Buktikan bahwa $A', B', C'$ bersifat
segaris bila dan hanya bila $\alpha\beta\gamma = 1$. \emph{(Tulislah
tiap titiknya sebagai \omterm{def:b2:affine:barycenter}{barisentrum} dua titik sudut; lalu tunjukkan bahwa tiga
titik segaris bila dan hanya bila baris koordinat barisentriknya terhadap
$(A, B, C)$ membentuk matriks $3 \times 3$ yang singular.)}
\end{exercise}

\begin{exercise}[$\star\star\star$]\label{exo:b2:affine:12}
Buktikan bahwa \omterm{def:b2:affine:convex}{bungkus cembung} himpunan bagian \omterm{def:b2:metric:compact}{kompak} $K$ milik $\R^n$
bersifat \omterm{def:b2:metric:compact}{kompak}. \emph{(Menurut \cref{thm:b2:affine:caratheodory},
$\operatorname{conv}(K)$ adalah peta sebuah himpunan \omterm{def:b2:metric:compact}{kompak} oleh
pemetaan yang \omterm{def:b2:metric:continuity}{kontinu}.)} Tunjukkanlah lewat sebuah contoh di $\R^2$ bahwa \omterm{def:b2:affine:convex}{bungkus
cembung} himpunan yang \emph{tertutup} tak harus tertutup.
\end{exercise}

\section{Soal: dari Radon ke Helly, titik penyeimbang dan teorema
Jung}

\begin{omfigure}
\begin{tikzpicture}[scale=0.9]
  \draw[omDef, thick] (0,0) -- (3,0) -- (0.9,2.6) -- cycle;
  \fill (0,0) circle (1.6pt) node[below left] {\small $A_1$};
  \fill (3,0) circle (1.6pt) node[below right] {\small $A_2$};
  \fill (0.9,2.6) circle (1.6pt) node[above] {\small $A_3$};
  \fill[omThm] (1.3,0.9) circle (1.8pt)
    node[below, omThm] {\small $A_4$};
\end{tikzpicture}
\qquad
\begin{tikzpicture}[scale=0.9]
  \draw[omDef, thick] (0,0) -- (3,0.3) -- (3.3,2.4) --
    (0.4,2.2) -- cycle;
  \draw[omThm, thick] (0,0) -- (3.3,2.4);
  \draw[omThm, thick] (3,0.3) -- (0.4,2.2);
  \fill (0,0) circle (1.6pt) node[below left] {\small $A_1$};
  \fill (3,0.3) circle (1.6pt) node[below right] {\small $A_2$};
  \fill (3.3,2.4) circle (1.6pt) node[above right] {\small $A_3$};
  \fill (0.4,2.2) circle (1.6pt) node[above left] {\small $A_4$};
  \fill[omProp] (1.75,1.27) circle (1.8pt);
\end{tikzpicture}

{\small Kedua tipe Radon bagi empat titik bidang pada
posisi umum: yakni satu titik di dalam segitiga yang lain
(dengan pemartisian $\{A_4\} \mid \{A_1, A_2, A_3\}$), atau posisi
\omterm{def:b2:affine:convex}{cembung}, tempat titik Radonnya (jingga) adalah perpotongan
kedua diagonalnya.}
\end{omfigure}

\begin{problem}[{Soal akhir pekan --- teorema Helly dan dua
panennya}]\label{pb:b2:affine:1}
Lema Radon (\cref{exo:b2:affine:7}) mengatakan bahwa $n + 2$
titik pada \omterm{def:b2:affine:def}{ruang afin} berdimensi $n$ selalu terpecah menjadi
dua kelompok yang \omterm{def:b2:affine:convex}{bungkus cembungnya} berpotongan. Soal ini mengubah
satu fakta aljabar linear itu menjadi rantai teorema geometri
kombinatorik: yakni teorema perpotongan Helly, teorema
titik penyeimbang (yakni sebuah median dua dimensi), dan teorema
peliputan Jung. Di sepanjang soal ini, bidangnya adalah $\R^2$ dengan
struktur Euklides biasanya, dan $\det$ adalah \omterm{def:b2:linalg:det}{determinan} dalam
basis kanoniknya.\index{teorema Helly}\index{lema
Radon}\index{teorema Jung}\index{titik penyeimbang}

\medskip
\noindent\textbf{Bagian I --- Koordinat barisentrik.} Titik
$A_0, \dots, A_k$ disebut \emph{bebas secara \omterm{def:b2:affine:subspace}{afin}} apabila
vektor $\vect{A_0A_1}, \dots, \vect{A_0A_k}$ bebas
secara linear.
\begin{enumerate}
  \item Tunjukkan bahwa kebebasan \omterm{def:b2:affine:subspace}{afinnya} tak bergantung pada
        pemilihan titik pangkalnya $A_0$, dan bahwa ia
        setara dengan: setiap kali dua keluarga bobot, yang masing-masing
        berjumlah $1$, mendefinisikan \omterm{def:b2:affine:barycenter}{barisentrum} $(A_0,
        \dots, A_k)$ yang sama, maka bobotnya berimpit.
  \item Misalkan $(A, B, C)$ bebas secara \omterm{def:b2:affine:subspace}{afin} pada bidangnya.
        Tunjukkan bahwa setiap titik $M$ menerima satu tripel tunggal
        $(\alpha, \beta, \gamma)$ dengan $\alpha + \beta +
        \gamma = 1$ dan $M = \operatorname{bar}(A, \alpha; B,
        \beta; C, \gamma)$ --- yakni \emph{koordinat
        barisentriknya}.
  \item Buktikan rumus \omterm{def:b2:linalg:det}{determinannya}
        \[
        \alpha = \frac{\det(\vect{MB},
        \vect{MC})}{\det(\vect{AB}, \vect{AC})},
        \qquad
        \beta = \frac{\det(\vect{MC},
        \vect{MA})}{\det(\vect{AB}, \vect{AC})},
        \qquad
        \gamma = \frac{\det(\vect{MA},
        \vect{MB})}{\det(\vect{AB}, \vect{AC})} :
        \]
        jadi koordinat barisentriknya adalah nisbah luas bertanda.
  \item Adapun garis $BC$, $CA$, $AB$ merupakan garis koordinat
        $\{\alpha = 0\}$, $\{\beta = 0\}$, $\{\gamma = 0\}$.
        Tunjukkan bahwa $M$ terletak di segitiga tertutup
        $\operatorname{conv}\{A, B, C\}$ bila dan hanya bila $\alpha, \beta,
        \gamma \geq 0$, dan bahwa ketiga garisnya memotong bidangnya
        menjadi tepat tujuh daerah, yang digolongkan tanda
        $(\alpha, \beta, \gamma)$ (dengan pola tanda $(-, -,
        -)$ yang mustahil).
  \item Misalkan $u \colon \R^2 \to \R$ \omterm{def:b2:affine:subspace}{pemetaan afin} (yakni sebuah
        \emph{bentuk \omterm{def:b2:affine:subspace}{afin}}). Tunjukkan $u(M) = \alpha\,u(A) +
        \beta\,u(B) + \gamma\,u(C)$, bahwa himpunan aras sebuah
        bentuk \omterm{def:b2:affine:subspace}{afin} yang tak tetap berupa garis, bahwa setiap garis
        muncul demikian, dan bahwa setengah bidang tertutup $\{u
        \geq c\}$ bersifat \omterm{def:b2:affine:convex}{cembung}.
\end{enumerate}

\medskip
\noindent\textbf{Bagian II --- Pemartisian Radon, diperhalus.} Sebuah
keluarga berisi $n + 2$ titik $\R^n$ disebut pada \emph{posisi
umum} apabila setiap $n + 1$ di antaranya bebas secara \omterm{def:b2:affine:subspace}{afin}.
Adapun sebuah \emph{kebergantungan \omterm{def:b2:affine:subspace}{afin}} bagi $(A_1, \dots, A_{n+2})$ adalah
keluarga $(\mu_i)$ dengan $\sum_i \mu_i = 0$ dan $\sum_i
\mu_i\,\vect{OA_i} = 0$ untuk satu (sehingga setiap) titik asal $O$.
\begin{enumerate}[resume]
  \item Hitunglah sebuah kebergantungan \omterm{def:b2:affine:subspace}{afin} tak nol bagi keempat
        titik $A_1 = (0,0)$, $A_2 = (3,0)$, $A_3 = (0,3)$,
        $A_4 = (1,1)$; lalu berikan pemartisian Radonnya beserta titik
        Radonnya.
  \item Tunjukkan bahwa untuk titik pada posisi umum, ruang
        vektor kebergantungan \omterm{def:b2:affine:subspace}{afinnya} berdimensi tepat
        $1$, dan bahwa kebergantungan tak nol \emph{tak} punya
        koefisien yang lenyap.
  \item Turunkan bahwa pemartisian Radon bagi $n + 2$ titik pada
        posisi umum bersifat tunggal (hingga pertukaran kedua
        bloknya), dengan tiap bloknya berupa himpunan indeks tempat
        $\mu_i$ bertanda tetap.
  \item Untuk empat titik bidang pada posisi umum,
        tunjukkanlah dikotominya: entah pemartisiannya bertipe
        $(1, 3)$ --- yakni satu titik di dalam segitiga
        ketiga lainnya --- atau bertipe $(2, 2)$: yakni keempat titiknya
        berada pada posisi \omterm{def:b2:affine:convex}{cembung} dan ruas yang menghubungkan
        kedua pasangannya (yakni diagonalnya) berpotongan, di titik
        Radonnya.
  \item Kerjakanlah pertanyaan 6 untuk bujur sangkar satuan $(0,0)$,
        $(1,0)$, $(1,1)$, $(0,1)$: yakni kebergantungannya, pemartisiannya, dan
        titik Radonnya.
\end{enumerate}

\medskip
\noindent\textbf{Bagian III --- Teorema Helly pada bidang.}
\begin{enumerate}[resume]
  \item Misalkan $C_1, C_2, C_3, C_4$ himpunan bagian \omterm{def:b2:affine:convex}{cembung} $\R^2$,
        yang setiap tiganya punya titik bersama. Pilihlah $x_i \in
        \bigcap_{j \neq i} C_j$ lalu terapkan lema Radon pada
        $x_1, \dots, x_4$: tunjukkanlah bahwa titik Radonnya termasuk
        keempat himpunannya. \emph{(Sebab untuk tiap $k$, blok yang tak
        memuat $x_k$ terdiri atas titik $C_k$.)}
  \item (Helly) Misalkan $C_1, \dots, C_m$ (dengan $m \geq 3$) himpunan bagian
        \omterm{def:b2:affine:convex}{cembung} $\R^2$, yang setiap tiganya berpotongan.
        Buktikan $\bigcap_{i=1}^m C_i \neq \emptyset$, secara
        induktif atas $m$: gantilah $C_{m-1}$ dan $C_m$ dengan
        $C_{m-1} \cap C_m$ lalu periksalah hipotesisnya bagi
        keluarga barunya lewat pertanyaan 11.
  \item Tiga contoh tandingan, satu bagi tiap hipotesisnya: (a) ketiga
        rusuk tertutup sebuah segitiga berpotongan berpasangan
        tetapi tak punya titik bersama (jadi $3$ tak dapat diturunkan menjadi
        $2$); (b) keempat himpunan $S_i = \{x_1, \dots, x_4\}
        \setminus \{x_i\}$, bagi empat titik pada posisi umum,
        memenuhi hipotesis perpotongan bertiganya
        tetapi tak memenuhi kesimpulannya (jadi kecembungannya berperan); (c)
        setengah bidang tertutup $H_k = \intco{k}{+\infty} \times
        \R$ dengan $k \in \N$ berpotongan berpasangan dan bertiga
        tetapi $\bigcap_k H_k = \emptyset$ (jadi keluarga tak hingga
        menuntut kekompakan).
  \item (Helly \omterm{def:b2:metric:compact}{kompak}) Misalkan $(K_i)_{i \in I}$ sembarang
        keluarga himpunan bagian \omterm{def:b2:metric:compact}{kompak} \omterm{def:b2:affine:convex}{cembung} $\R^2$, yang setiap tiganya
        berpotongan. Dengan memakai pertanyaan 12 dan sifat
        Borel--Lebesgue
        (\cref{thm:b2:metric:borellebesgue}), tunjukkanlah
        $\bigcap_{i \in I} K_i \neq \emptyset$.
  \item (Panen pertama) Misalkan $S$ himpunan berhingga titik pada
        bidang dan $r > 0$. Tunjukkan: bahwa bila setiap tiga titik
        $S$ terletak pada suatu cakram tertutup berjari-jari $r$, maka $S$
        terletak pada satu cakram tertutup berjari-jari $r$. \emph{(Terapkan
        Helly pada cakram $\overline D(p, r)$ dengan $p \in S$.)}
\end{enumerate}

\medskip
\noindent\textbf{Bagian IV --- Teorema titik penyeimbang.} Sebuah
\emph{titik penyeimbang} bagi himpunan berhingga $S$ berisi $n$ titik pada
bidang adalah titik $c$ (yang tak harus di $S$) sedemikian sehingga setiap
setengah bidang tertutup yang memuat $c$ memuat sekurangnya $n/3$
titik $S$.
\begin{enumerate}[resume]
  \item (Dimensi $1$) Untuk bilangan real $x_1 \leq \dots \leq x_n$,
        tunjukkan bahwa mediannya $c = x_{\lceil n/2 \rceil}$
        memenuhi: bahwa setiap setengah garis tertutup yang memuat $c$
        memuat sekurangnya $n/2$ di antara $x_i$.
  \item (Lema pencacahan) Bila $A, B, C$ himpunan bagian $S$ dengan
        $\abs A, \abs B, \abs C > \tfrac{2n}3$, tunjukkanlah $A \cap
        B \cap C \neq \emptyset$.
  \item Misalkan $m = \floor{2n/3} + 1$ dan misalkan $\mathcal F$
        keluarga (yang berhingga) berisi \omterm{def:b2:affine:convex}{bungkus cembung}
        $\operatorname{conv}(T)$ dengan $T \subseteq S$ dan $\abs T =
        m$. Tunjukkan bahwa setiap tiga anggota $\mathcal F$ punya
        titik bersama, lalu turunkan dari Helly sebuah titik $c$ yang
        bersama bagi semuanya.
  \item Buktikan bahwa $c$ ini merupakan titik penyeimbang $S$: itulah
        \emph{teorema titik penyeimbang}. \emph{(Sebab bila sebuah setengah
        bidang tertutup lewat $c$ memuat kurang dari $n/3$
        titik, maka komplemen \omterm{def:b2:metric:topology}{terbukanya} akan memuat sebuah himpunan $T$
        berisi $m$ titik, sehingga $\operatorname{conv}(T)$ akan
        menghindari $c$.)}
  \item Ketajamannya: misalkan $n = 3k$ lalu tempatkanlah $k$ titik pada
        masing-masing tiga cakram berjari-jari kecil $\varepsilon$ yang berpusat
        di titik sudut sebuah segitiga besar. Tunjukkan bahwa untuk
        \emph{setiap} titik $c$ pada bidangnya ada setengah
        bidang tertutup yang memuat $c$ dan memuat paling banyak $n/3$
        titik $S$, jadi konstanta $1/3$ tak dapat
        diperbaiki. \emph{(Sebab di antara ketiga arah dari $c$
        ke pusat cakramnya, dua di antaranya bersudut paling banyak
        $2\pi/3$.)}
\end{enumerate}

\medskip
\noindent\textbf{Bagian V --- Teorema Jung dan rangkuman.}
\begin{enumerate}[resume]
  \item (Lema segitiga) Misalkan $P, Q, R$ tiga titik dengan
        jarak berpasangan $\leq 1$. Tunjukkan bahwa ketiganya terletak pada
        cakram tertutup berjari-jari $1/\sqrt3$. \emph{(Bila suatu
        sudutnya $\geq \pi/2$, ambillah cakram bergaris tengah sisi
        terpanjangnya, dengan memakai rumus garis berat
        $\norm{RM}^2 = \tfrac12\norm{RP}^2 +
        \tfrac12\norm{RQ}^2 - \tfrac14\norm{PQ}^2$; sedangkan bila
        segitiganya lancip, batasilah jari-jari lingkar luarnya $a/(2\sin
        \widehat A)$ lewat sudut terbesarnya, yang terletak di
        $\intco{\pi/3}{\pi/2}$.)}
  \item (Jung) Turunkan: bahwa setiap himpunan bagian \omterm{def:b2:metric:compact}{kompak} bidang yang
        berdiameter $\leq 1$ termuat pada sebuah cakram tertutup
        berjari-jari $1/\sqrt3$.
  \item Ketajamannya: untuk segitiga sama sisi $A_1A_2A_3$
        bersisi $1$ dengan titik berat $G$, buktikanlah kesamaan
        Leibniz $\sum_i \norm{\vect{OA_i}}^2 =
        3\norm{\vect{OG}}^2 + \sum_i \norm{\vect{GA_i}}^2$
        untuk setiap titik $O$, lalu simpulkan bahwa sembarang cakram yang
        memuat ketiga titik sudutnya berjari-jari $\geq
        1/\sqrt3$, dengan kesamaan hanya bagi cakram lingkar luarnya.
  \item (Helly di $\R^n$) Nyatakan lalu buktikan teorema Helly di
        $\R^n$: bahwa bila berhingga banyak \omterm{def:b2:affine:convex}{himpunan cembung} sedemikian sehingga
        setiap $n + 1$ di antaranya berpotongan, maka semuanya
        berpotongan. \emph{(Lema Radon
        \cref{exo:b2:affine:7} menangani $n + 2$ himpunan; lalu
        berinduksilah seperti pada pertanyaan 12.)}
  \item Rangkuman. Rangkailah rantainya
        \[
        \text{kebergantungan afin} \Rightarrow \text{Radon}
        \Rightarrow \text{Helly} \Rightarrow
        \text{titik penyeimbang dan Jung},
        \]
        dengan menunjukkan dalam satu kalimat masing-masing: di mana aljabar linear
        masuk, di mana tanda bobotnya masuk, di mana
        kecembungannya masuk, dan langkah tunggal mana yang memakai
        dimensi bidangnya. Lalu konstanta $3$ (pada
        Helly), $1/3$ (titik penyeimbang) dan $1/\sqrt3$ (Jung)
        menjadi apa di $\R^n$? (Nyatakan saja tanpa bukti.)

\end{enumerate}
\end{problem}
